\begin{textsample}{POS Dim 5 – global\_north\_2025\_01 – Score 16.00 – global\_north\_2025\_01/120719318.txt}  \label{ex:f5_pos_268}
A judge has ruled that a \textbf{legal} challenge to UK arms \textbf{exports} to Israel can move forward in the High \textbf{Court}, with the \textbf{case} set to focus on the government ' s decision to continue sending F-35 fighter jet components that could end up in Israeli planes.
In a judgment filed on Thursday, Mr Justice Chamberlain said he was granting permission for a judicial review that will look specifically at grounds that challenge the UK ' s"F-35 Carve Out".
If the challenge is successful, the government ' s decision to \textbf{license} the \textbf{export} of F-35 components"will have to be revisited", the judge wrote.
While a multi-phase ceasefire deal brought a cessation of hostilities in Gaza earlier this month, the judge said he had proceeded on the basis that the ceasefire"makes no difference to the substantive or procedural issues in this \textbf{case}".
This is the latest development in the challenge first brought by Palestinian human rights organisation Al-Haq and the UK-based Global \textbf{Legal} Action Network ( Glan ) against the UK @ @ @ @ @ @ @ @ @ @ in October 2023.
New MEE newsletter : Jerusalem Dispatch Sign up to get the latest \textbf{insights} and \textbf{analysis} on Israel-Palestine, alongside Turkey Unpacked and other MEE newsletters At least 47,283 Palestinians were killed in the 15-month war according to the latest figures from the health ministry in Gaza.
After the \textbf{ruling}, Shawan Jabarin, general-director of Al-Haq, said :"Gaza is destroyed, it is unliveable.
Palestinians in Gaza have been killed and erased by weapons whose components are supplied to Israel by the UK Government, \textbf{acting} in full knowledge of the consequences."Jennine Walker, a lawyer with Glan, said :"Finally, over a year after this \textbf{case} was started, a \textbf{court} will decide Al Haq and Glan ' s \textbf{legal} challenge."The organisations ' request for a judicial review was initially rejected, but then accepted on appeal last April.
At that time, the UK government ' s position was that there was no clear risk that arms \textbf{exported} from the UK might be used @ @ @ @ @ @ @ @ @ @ in Gaza.
On three separate occasions - on 18 December 2023, 8 April and 28 May 2024 - the government decided to proceed with \textbf{licensing} all arms \textbf{exports} to Israel, according to \textbf{court} documents.
War \textbf{crime} risk This changed last September when the new Labour government suspended 30 arms \textbf{export} licences after a review found there was indeed a clear risk that Israel might use UK arms to \textbf{commit} war \textbf{crimes} in Gaza.
But the UK continued to \textbf{export} the fighter jet parts to a global pool that could end up in Israeli F-35s, which experts have said were critical to its air campaign against Gaza.
The government has acknowledged that there is a clear risk that Israel may \textbf{commit} war \textbf{crimes} using F-35s, but has said it could not suspend the parts without disrupting the entire F-35 fleet and threatening global peace and security"within weeks".
New report lays out full extent of UK-Israel military partnership in Gaza Campaigners argue that while the government assessed the risks of suspending the F-35 @ @ @ @ @ @ @ @ @ @ to Palestinians and the \textbf{international} rule of law of continuing to \textbf{exports} the parts, and that it has failed to follow its own guidelines and \textbf{international} treaty obligations.
In addition to the focus on the F-35s, the \textbf{case} moving forward will also look at the UK government ' s decision to continue supplying equipment that the government assessed that Israel would not be capable of using in Gaza, despite finding that Israel was not \textbf{committed} to complying with \textbf{international} humanitarian law.
Oxfam, \textbf{Amnesty} \textbf{International} and Human Rights Watch are interveners in the \textbf{case}, providing evidence and testimonies to support Glan and Al-Haq ' s claim.
Helen Stawski, Oxfam ' s head of policy, said she welcomed the \textbf{ruling}, but said that given the government has"acknowledged that Israel is breaching \textbf{international} humanitarian law in Gaza, it must immediately suspend all arms sales to Israel, rather than waiting on a \textbf{court} ' s judgement".
She added :"Leaving a loophole for F35 components is unconscionable and out of step with the Government ' s domestic @ @ @ @ @ @ @ @ @ @ to meet even with the current pause in hostilities."A hearing that is expected to last at least three days has been provisionally set for 12 May.
Middle East Eye delivers \textbf{independent} and unrivalled coverage and \textbf{analysis} of the Middle East, North Africa and beyond.
To learn more about republishing this content and the \textbf{associated} \textbf{fees}, please fill out this form.
More about MEE can be found here.

% matched lemmas: act, amnesty, analysis, associate, case, commit, court, crime, export, fee, independent, insight, international, legal, license, ruling
\end{textsample}
